\section{Phase 1: Reconstructing the Global MINCO Mother Problem as the Reference Model for BLOM}
\label{sec:phase1_minco_reference}

In this section we reconstruct the mathematical backbone of MINCO under the notation fixed in
Section~\ref{sec:blom_phase0}. The purpose is not to reproduce the full GCOPTER planning
framework, but to isolate the precise chain
\[
  \text{mother problem}
  \;\Longrightarrow\;
  \text{optimality \& uniqueness}
  \;\Longrightarrow\;
  \text{banded coefficient system}
  \;\Longrightarrow\;
  \text{parameter sensitivity \& gradients},
\]
which will later be localized in the BLOM construction.

Throughout this section we remain in the canonical scalar setting of Phase~0:
\[
  p:[0,\mathcal T]\to \mathbb R,
  \qquad
  \mathcal T=\sum_{i=1}^{M}T_i,
  \qquad
  T_i>0,
  \qquad
  s=4,
  \qquad
  d_i=1.
\]
Hence the control effort is the snap energy
\[
  J[p]
  =
  \sum_{i=1}^{M}\int_{0}^{T_i}\bigl|p_i^{(4)}(t)\bigr|^2\,dt.
\]
The sparse external parameter vector is still
\begin{equation}
  \label{eq:phase1_theta}
  \theta=(q_1,\dots,q_{M-1},T_1,\dots,T_M)^\top,
\end{equation}
exactly as fixed in Phase~0. The boundary waypoints $q_0$ and $q_M$ remain fixed data.

\subsection{Why the MINCO mother problem must strengthen the Phase~0 reference problem}
\label{subsec:phase1_why_strengthen}

The global reference problem in Phase~0 is sufficient as a \emph{comparison target} for BLOM,
but it is not yet a mathematically adequate MINCO mother problem for uniqueness. Indeed, if
one fixes only waypoint values and minimizes the snap energy, then uniqueness generally fails.
For example, when $M=1$, any cubic interpolant between $q_0$ and $q_1$ has zero snap cost.
Therefore, in the MINCO mother problem we must additionally fix the full boundary jet up to
order $s-1=3$ at both ends.

Accordingly, we introduce fixed boundary data
\begin{equation}
  \label{eq:phase1_boundary_jets}
  \zeta^-=(\zeta^-_0,\zeta^-_1,\zeta^-_2,\zeta^-_3)^\top,
  \qquad
  \zeta^+=(\zeta^+_0,\zeta^+_1,\zeta^+_2,\zeta^+_3)^\top,
\end{equation}
with
\[
  \zeta^-_0=q_0,
  \qquad
  \zeta^+_0=q_M.
\]
In the canonical minimum-snap setting, these are the position, velocity, acceleration, and jerk
values at the initial and terminal times.

\subsection{Mother problem}
\label{subsec:phase1_mother_problem}

Let
\[
  t_0:=0,
  \qquad
  t_i:=\sum_{r=1}^{i}T_r,
  \qquad
  i=1,\dots,M,
\]
so that $t_M=\mathcal T$.
We consider trajectories $p$ whose restriction to each interval $[t_{i-1},t_i]$ is denoted by
\[
  p_i:[0,T_i]\to\mathbb R,
  \qquad
  p_i(\tau)=p(t_{i-1}+\tau).
\]

Define the admissible set
\begin{equation}
  \label{eq:phase1_admissible_set}
  \mathcal A(q,T;\zeta^-,\zeta^+)
  =
  \left\{
    p
    \;\middle|\;
    \begin{aligned}
      & p_i\in H^4(0,T_i), \quad i=1,\dots,M,\\
      & p(t_i)=q_i,\quad i=0,\dots,M,\\
      & p^{(\ell)}(0)=\zeta^-_\ell,\quad p^{(\ell)}(\mathcal T)=\zeta^+_\ell,
      \quad \ell=1,2,3,\\
      & p^{(\ell)}(t_i^-)=p^{(\ell)}(t_i^+),\quad \ell=1,2,3,\quad i=1,\dots,M-1
    \end{aligned}
  \right\}.
\end{equation}
The Phase~1 MINCO mother problem is
\begin{equation}
  \label{eq:phase1_mother_problem}
  \boxed{
    \min_{p\in \mathcal A(q,T;\zeta^-,\zeta^+)}
    \;
    J[p]
    =
    \sum_{i=1}^{M}\int_{0}^{T_i}\bigl|p_i^{(4)}(\tau)\bigr|^2\,d\tau.
  }
\end{equation}

\paragraph{Remark on the role of \texorpdfstring{$d_i=1$}{d_i=1}.}
The canonical choice $d_i=1$ means that only the waypoint value $p(t_i)=q_i$ is prescribed at each
interior knot. No interior derivative value is fixed. The continuity of orders $1,2,3$ belongs to the
integrator-chain state continuity, while the continuity of higher orders will emerge from optimality.

\subsection{Optimality conditions and uniqueness}
\label{subsec:phase1_optimality_uniqueness}

We now state and prove the canonical mother theorem for the Phase~0 setting $(s,d_i)=(4,1)$.

\begin{theorem}[Optimality conditions and uniqueness in the canonical minimum-snap case]
  \label{thm:phase1_minco_theorem}
  Fix $T_i>0$ for all $i$, waypoint data $q_0,\dots,q_M\in\mathbb R$, and boundary jets
  $\zeta^-,\zeta^+\in\mathbb R^4$ with $\zeta^-_0=q_0$ and $\zeta^+_0=q_M$.
  A trajectory $p^\star\in\mathcal A(q,T;\zeta^-,\zeta^+)$ solves
  \eqref{eq:phase1_mother_problem} if and only if the following hold:

  \begin{enumerate}
    \item For each $i=1,\dots,M$, the segment $p_i^\star$ is a polynomial of degree at most $7$.

    \item The boundary conditions are satisfied:
      \[
        (p^\star)^{(\ell)}(0)=\zeta^-_\ell,
        \qquad
        (p^\star)^{(\ell)}(\mathcal T)=\zeta^+_\ell,
        \qquad \ell=0,1,2,3,
      \]
      where $\zeta^-_0=q_0$ and $\zeta^+_0=q_M$.

    \item The waypoint conditions are satisfied:
      \[
        p^\star(t_i)=q_i,\qquad i=1,\dots,M-1.
      \]

    \item The trajectory is $C^6$ at every interior knot:
      \[
        (p^\star)^{(\ell)}(t_i^-)
        =
        (p^\star)^{(\ell)}(t_i^+),
        \qquad
        \ell=0,1,\dots,6,
        \qquad
        i=1,\dots,M-1.
      \]
  \end{enumerate}

  Moreover, the solution of \eqref{eq:phase1_mother_problem} exists and is unique.
\end{theorem}

\begin{proof}
  We split the proof into necessity, sufficiency, and uniqueness.

  \paragraph{Necessity.}
  Let $p^\star\in\mathcal A(q,T;\zeta^-,\zeta^+)$ be optimal and let
  $h\in \mathcal T_{p^\star}\mathcal A$ be any admissible variation. Then
  \[
    h(t_i)=0,\qquad i=0,\dots,M,
  \]
  and
  \[
    h^{(\ell)}(0)=h^{(\ell)}(\mathcal T)=0,\qquad \ell=1,2,3,
  \]
  while
  \[
    h^{(\ell)}(t_i^-)=h^{(\ell)}(t_i^+),\qquad \ell=1,2,3,\quad i=1,\dots,M-1.
  \]
  Since $J$ is quadratic,
  \[
    0
    =
    \frac{1}{2}\,\delta J[p^\star;h]
    =
    \sum_{i=1}^{M}\int_{0}^{T_i} (p_i^\star)^{(4)}(\tau)\, h_i^{(4)}(\tau)\,d\tau.
  \]
  Integrating by parts four times on each segment gives
  \begin{align}
    \int_{0}^{T_i} (p_i^\star)^{(4)} h_i^{(4)}\,d\tau
    &=
    \Bigl[
      (p_i^\star)^{(4)} h_i^{(3)}
      -
      (p_i^\star)^{(5)} h_i^{(2)}
      +
      (p_i^\star)^{(6)} h_i^{(1)}
      -
      (p_i^\star)^{(7)} h_i
    \Bigr]_{0}^{T_i}
    \notag\\
    &\qquad
    +
    \int_{0}^{T_i}(p_i^\star)^{(8)}(\tau)\,h_i(\tau)\,d\tau .
    \label{eq:phase1_ibp4}
  \end{align}
  Summing \eqref{eq:phase1_ibp4} over $i$, the endpoint contributions at $t=0$ and $t=\mathcal T$
  vanish because the boundary jet is fixed, and the terms multiplying $h(t_i)$ vanish because
  $h(t_i)=0$. Hence
  \begin{align}
    0
    &=
    \sum_{i=1}^{M}\int_{0}^{T_i}(p_i^\star)^{(8)}(\tau)\,h_i(\tau)\,d\tau
    \notag\\
    &\quad
    +
    \sum_{i=1}^{M-1}
    \Bigl[
      \bigl((p_i^\star)^{(4)}(T_i)-(p_{i+1}^\star)^{(4)}(0)\bigr)\,h^{(3)}(t_i)
      \notag\\
      &\qquad\qquad
      -
      \bigl((p_i^\star)^{(5)}(T_i)-(p_{i+1}^\star)^{(5)}(0)\bigr)\,h^{(2)}(t_i)
      \notag\\
      &\qquad\qquad
      +
      \bigl((p_i^\star)^{(6)}(T_i)-(p_{i+1}^\star)^{(6)}(0)\bigr)\,h^{(1)}(t_i)
    \Bigr].
    \label{eq:phase1_variation_reduced}
  \end{align}
  Now $h$ is arbitrary in the interior of each segment, and the traces
  $h^{(1)}(t_i),h^{(2)}(t_i),h^{(3)}(t_i)$ are independent free variables. Therefore
  \[
    (p_i^\star)^{(8)}(\tau)=0
    \quad\text{for all }\tau\in(0,T_i),
  \]
  hence each $p_i^\star$ is a polynomial of degree at most $7$, and
  \[
    (p_i^\star)^{(\ell)}(T_i)=(p_{i+1}^\star)^{(\ell)}(0),
    \qquad \ell=4,5,6,
    \qquad i=1,\dots,M-1.
  \]
  Since $p^\star\in\mathcal A(q,T;\zeta^-,\zeta^+)$ already has continuity of orders $0,1,2,3$,
  we conclude that $p^\star$ is $C^6$ at every interior knot.

  \paragraph{Sufficiency.}
  Assume a feasible trajectory $p$ satisfies items (1)--(4). Then each segment is a polynomial of
  degree at most $7$, so $p_i^{(8)}\equiv 0$ on every segment. Repeating the integration-by-parts
  identity \eqref{eq:phase1_ibp4}, all knot terms in \eqref{eq:phase1_variation_reduced} vanish
  because of $C^6$ continuity and $h(t_i)=0$. Thus
  \[
    \delta J[p;h]=0
    \qquad
    \text{for all } h\in \mathcal T_p\mathcal A .
  \]
  Since $J$ is a quadratic functional with nonnegative Hessian,
  \[
    \delta^2 J[p;h]
    =
    2\sum_{i=1}^{M}\int_{0}^{T_i}|h_i^{(4)}(\tau)|^2\,d\tau
    \ge 0,
  \]
  the vanishing of the first variation implies that $p$ is a global minimizer over the affine feasible
  set.

  \paragraph{Uniqueness.}
  Suppose $p^\star$ and $\tilde p^\star$ are two minimizers. Their difference
  $h:=p^\star-\tilde p^\star$ satisfies all homogeneous constraints, and
  \[
    0
    =
    J[p^\star]+J[\tilde p^\star]-2J\!\left[\frac{p^\star+\tilde p^\star}{2}\right]
    =
    \frac{1}{2}\sum_{i=1}^{M}\int_{0}^{T_i}|h_i^{(4)}(\tau)|^2\,d\tau.
  \]
  Hence $h_i^{(4)}\equiv 0$ for every $i$, so each $h_i$ is a polynomial of degree at most $3$.
  Because $h$ is $C^3$ across all interior knots, $h$ is in fact a single global cubic polynomial on
  $[0,\mathcal T]$. Moreover, all boundary jet data are homogeneous:
  \[
    h^{(\ell)}(0)=0,\qquad \ell=0,1,2,3.
  \]
  Therefore the global cubic polynomial $h$ has all four Taylor coefficients at $0$ equal to zero,
  which implies $h\equiv 0$. Thus $p^\star=\tilde p^\star$, and the minimizer is unique.
\end{proof}

\begin{remark}[Why the endpoint jet is not optional]
  \label{rem:phase1_endpoint_jet}
  The fixed endpoint jet in \eqref{eq:phase1_boundary_jets} is not a cosmetic addition but a
  mathematical necessity for uniqueness. If one keeps only the waypoint constraints
  $p(t_i)=q_i$ and removes the higher-order endpoint conditions, then the functional
  \eqref{eq:phase1_mother_problem} is no longer strictly convex on the feasible affine space.
  For example, when $M=1$, every cubic interpolant between $q_0$ and $q_1$ has zero snap energy.
\end{remark}

\subsection{Banded linear system for the optimal coefficient vector}
\label{subsec:phase1_banded_system}

Theorem~\ref{thm:phase1_minco_theorem} converts the variational problem
\eqref{eq:phase1_mother_problem} into a purely algebraic construction.

For each segment we adopt the monomial basis
\[
  \beta(\tau)
  =
  (1,\tau,\tau^2,\tau^3,\tau^4,\tau^5,\tau^6,\tau^7)^\top \in \mathbb R^8,
\]
and write
\begin{equation}
  \label{eq:phase1_segment_poly}
  p_i(\tau)=c_i^\top \beta(\tau),
  \qquad
  \tau\in[0,T_i],
  \qquad
  c_i\in\mathbb R^8.
\end{equation}
Let
\[
  c=(c_1^\top,\dots,c_M^\top)^\top \in \mathbb R^{8M}.
\]
For convenience define the derivative row operators
\[
  B^{(\ell)}(\tau)
  :=
  \frac{d^\ell}{d\tau^\ell}\beta(\tau)^\top
  \in \mathbb R^{1\times 8},
  \qquad
  \ell=0,1,\dots,7.
\]
Then
\[
  p_i^{(\ell)}(\tau)=B^{(\ell)}(\tau)c_i.
\]

\paragraph{Boundary rows.}
Define
\begin{equation}
  \label{eq:phase1_boundary_blocks}
  F_0
  =
  \begin{bmatrix}
    B^{(0)}(0)\\
    B^{(1)}(0)\\
    B^{(2)}(0)\\
    B^{(3)}(0)
  \end{bmatrix}
  \in\mathbb R^{4\times 8},
  \qquad
  E_M(T_M)
  =
  \begin{bmatrix}
    B^{(0)}(T_M)\\
    B^{(1)}(T_M)\\
    B^{(2)}(T_M)\\
    B^{(3)}(T_M)
  \end{bmatrix}
  \in\mathbb R^{4\times 8},
\end{equation}
and the boundary data vectors
\begin{equation}
  \label{eq:phase1_boundary_rhs}
  D_0=\zeta^- \in\mathbb R^4,
  \qquad
  D_M=\zeta^+ \in\mathbb R^4.
\end{equation}

\paragraph{Interior rows.}
For each interior knot $t_i$ with $1\le i\le M-1$, define
\begin{equation}
  \label{eq:phase1_interior_blocks}
  E_i(T_i)
  =
  \begin{bmatrix}
    B^{(0)}(T_i)\\
    B^{(0)}(T_i)\\
    B^{(1)}(T_i)\\
    B^{(2)}(T_i)\\
    B^{(3)}(T_i)\\
    B^{(4)}(T_i)\\
    B^{(5)}(T_i)\\
    B^{(6)}(T_i)
  \end{bmatrix}
  \in\mathbb R^{8\times 8},
  \qquad
  F_i
  =
  \begin{bmatrix}
    0_{1\times 8}\\
    -\,B^{(0)}(0)\\
    -\,B^{(1)}(0)\\
    -\,B^{(2)}(0)\\
    -\,B^{(3)}(0)\\
    -\,B^{(4)}(0)\\
    -\,B^{(5)}(0)\\
    -\,B^{(6)}(0)
  \end{bmatrix}
  \in\mathbb R^{8\times 8},
\end{equation}
and
\begin{equation}
  \label{eq:phase1_interior_rhs}
  d_i
  =
  (q_i,0,0,0,0,0,0,0)^\top \in\mathbb R^8.
\end{equation}
The first row of \eqref{eq:phase1_interior_blocks} enforces the waypoint condition
$p_i(T_i)=q_i$, while the remaining seven rows enforce continuity of derivatives of orders
$0,\dots,6$ between segments $i$ and $i+1$.

Collecting all rows yields the linear system
\begin{equation}
  \label{eq:phase1_Mcb}
  M(T)c(q,T)=b(q),
\end{equation}
where
\begin{equation}
  \label{eq:phase1_M_matrix}
  M(T)
  =
  \begin{bmatrix}
    F_0        & 0          & 0          & \cdots & 0 \\
    E_1(T_1)   & F_1        & 0          & \cdots & 0 \\
    0          & E_2(T_2)   & F_2        & \ddots & \vdots \\
    \vdots     & \ddots     & \ddots     & \ddots & 0 \\
    0          & \cdots     & 0          & E_{M-1}(T_{M-1}) & F_{M-1} \\
    0          & \cdots     & 0          & 0      & E_M(T_M)
  \end{bmatrix}
  \in\mathbb R^{8M\times 8M},
\end{equation}
and
\begin{equation}
  \label{eq:phase1_b_vector}
  b(q)
  =
  \begin{bmatrix}
    D_0\\
    d_1\\
    d_2\\
    \vdots\\
    d_{M-1}\\
    D_M
  \end{bmatrix}
  \in\mathbb R^{8M}.
\end{equation}

\begin{proposition}[Structure of the coefficient system]
  \label{prop:phase1_banded}
  For fixed $s=4$, the matrix $M(T)$ in \eqref{eq:phase1_M_matrix} is block bi-diagonal with
  constant block size $8\times 8$. Consequently:
  \begin{enumerate}
    \item $M(T)$ has uniformly bounded bandwidth independent of $M$;
    \item the unique coefficient vector $c(q,T)$ is obtained by solving a banded linear system;
    \item a banded PLU factorization computes $c(q,T)$ in $O(M)$ time and $O(M)$ memory.
  \end{enumerate}
\end{proposition}

\begin{proof}
  The block pattern in \eqref{eq:phase1_M_matrix} is immediate from the fact that every interior
  optimality condition couples only two adjacent segments. Because the block size is fixed
  ($8\times 8$ for $s=4$), both the scalar bandwidth and the fill pattern remain $O(1)$ with respect
  to $M$. Theorem~\ref{thm:phase1_minco_theorem} implies existence and uniqueness of the
  trajectory, hence the coefficient vector is unique and $M(T)$ is nonsingular. Standard banded
  PLU factorization for nonsingular banded matrices therefore yields the claimed linear
  complexity.
\end{proof}

\paragraph{Canonical specialization to Phase~0.}
In the notation of Phase~0, the external sparse parameter vector remains
\[
  \theta=(q_1,\dots,q_{M-1},T_1,\dots,T_M)^\top.
\]
The only additional data in Phase~1 are the fixed boundary jets $\zeta^-$ and $\zeta^+$, which are
not optimization variables.

\subsection{Parameter sensitivity and linear-complexity gradient propagation}
\label{subsec:phase1_sensitivity_gradient}

The role of MINCO as a trajectory class is not merely to generate a reference trajectory, but to
support further spatial--temporal deformation under user-defined objectives. Following the
MINCO construction, let
\begin{equation}
  \label{eq:phase1_W_def}
  W(q,T)=K(c(q,T),T),
\end{equation}
where $K:\mathbb R^{8M}\times\mathbb R^M_{>0}\to\mathbb R$ is any $C^2$ objective (or smooth penalty)
with available partial derivatives $\partial K/\partial c$ and $\partial K/\partial T$.

Differentiate \eqref{eq:phase1_Mcb}:
\begin{equation}
  \label{eq:phase1_diff_Mcb}
  M\,dc + (dM)\,c = db.
\end{equation}
Hence
\begin{equation}
  \label{eq:phase1_dc_formula}
  dc
  =
  M^{-1}\bigl(db-(dM)c\bigr).
\end{equation}
Using \eqref{eq:phase1_W_def}, the total differential of $W$ is
\begin{equation}
  \label{eq:phase1_dW}
  dW
  =
  \left(\frac{\partial K}{\partial c}\right)^\top dc
  +
  \left(\frac{\partial K}{\partial T}\right)^\top dT.
\end{equation}
To eliminate the inverse of $M$, define the adjoint vector $g\in\mathbb R^{8M}$ by
\begin{equation}
  \label{eq:phase1_adjoint}
  M(T)^\top g = \frac{\partial K}{\partial c}.
\end{equation}
Substituting \eqref{eq:phase1_dc_formula} into \eqref{eq:phase1_dW} and using
\eqref{eq:phase1_adjoint} gives
\begin{equation}
  \label{eq:phase1_dW_compact}
  dW
  =
  g^\top db
  -
  g^\top (dM)c
  +
  \left(\frac{\partial K}{\partial T}\right)^\top dT.
\end{equation}

\paragraph{Gradient with respect to waypoint parameters.}
Only the vector $b(q)$ depends on $q$, and for $1\le i\le M-1$,
\[
  \frac{\partial b}{\partial q_i}
  =
  e_{4+8(i-1)+1},
\]
where $e_j$ is the $j$-th Euclidean basis vector in $\mathbb R^{8M}$. Therefore,
\begin{equation}
  \label{eq:phase1_grad_qi}
  \frac{\partial W}{\partial q_i}
  =
  g^\top \frac{\partial b}{\partial q_i}
  =
  g_{\,4+8(i-1)+1},
  \qquad
  i=1,\dots,M-1.
\end{equation}
Hence all spatial gradients are obtained once the single adjoint solve \eqref{eq:phase1_adjoint}
is completed.

\paragraph{Gradient with respect to segment durations.}
Only $M(T)$ depends on $T$. Therefore,
\begin{equation}
  \label{eq:phase1_grad_Ti_general}
  \frac{\partial W}{\partial T_i}
  =
  \frac{\partial K}{\partial T_i}
  -
  g^\top
  \left(\frac{\partial M}{\partial T_i}\right)c,
  \qquad
  i=1,\dots,M.
\end{equation}
Because $T_i$ appears only in the local right-end evaluation of segment $i$, the matrix
$\partial M/\partial T_i$ has support in exactly one block row:
\begin{equation}
  \label{eq:phase1_dM_support}
  \frac{\partial M}{\partial T_i}
  =
  \begin{cases}
    \text{zero everywhere except the block } \dfrac{\partial E_i(T_i)}{\partial T_i}, & 1\le i\le M-1,\\[1.2ex]
    \text{zero everywhere except the block } \dfrac{\partial E_M(T_M)}{\partial T_M}, & i=M.
  \end{cases}
\end{equation}
Thus each $\partial W/\partial T_i$ is computed from a constant-size local contraction, and the
whole vector $\partial W/\partial T$ is obtained in $O(M)$ total complexity.

\begin{proposition}[Linear-complexity gradient propagation]
  \label{prop:phase1_linear_gradient}
  Assume that $K(c,T)$ and its partial derivatives with respect to $c$ and $T$ are available.
  Then the gradients
  \[
    \frac{\partial W}{\partial q}
    \in \mathbb R^{M-1},
    \qquad
    \frac{\partial W}{\partial T}
    \in \mathbb R^{M}
  \]
  can be computed in $O(M)$ time and $O(M)$ memory.
\end{proposition}

\begin{proof}
  The forward coefficients $c(q,T)$ are obtained by one banded solve of \eqref{eq:phase1_Mcb},
  which is $O(M)$ by Proposition~\ref{prop:phase1_banded}. The adjoint vector $g$ is obtained by
  one banded solve of \eqref{eq:phase1_adjoint}, again in $O(M)$. Formula
  \eqref{eq:phase1_grad_qi} extracts one entry of $g$ for each $q_i$, and
  \eqref{eq:phase1_grad_Ti_general}--\eqref{eq:phase1_dM_support} require only constant-size local
  contractions for each $T_i$. Therefore the whole gradient propagation is linear in $M$.
\end{proof}

\subsection{Why Phase~1 is logically prior to BLOM}
\label{subsec:phase1_prior_to_blom}

The purpose of Phase~1 is not to advocate global MINCO as the final representation, but to fix
the exact mathematical chain that BLOM must later localize:
\[
  \boxed{
    \begin{aligned}
      &\text{variational mother problem}
      &&\Longrightarrow\text{optimality \& uniqueness}\\
      &\Longrightarrow M(T)c(q,T)=b(q)
      &&\Longrightarrow W(q,T)=K(c(q,T),T)\\
      &\Longrightarrow\left(\frac{\partial W}{\partial q},\frac{\partial W}{\partial T}\right)
    \end{aligned}
  }
\]
All later BLOM theorems will be local modifications of one of these five arrows, not independent
constructions detached from them.

\begin{remark}[General-\texorpdfstring{$s$}{s} version]
  \label{rem:phase1_general_s}
  The whole section extends verbatim to arbitrary control order $s\ge 2$.
  The replacements are
  \[
    4 \mapsto s,
    \qquad
    7 \mapsto 2s-1,
    \qquad
    6 \mapsto 2s-2,
    \qquad
    8 \mapsto 2s.
  \]
  For an interior knot with prescribed derivative count $d_i\le s$, the optimality theorem gives
  continuity order $2s-d_i-1$. In the canonical Phase~0 setting, $d_i=1$, hence the interior
  continuity becomes $2s-2$, which is $6$ when $s=4$.
\end{remark}